\section{Extension: A Predictor Correlated with Fundamentals}\label{sec:correlated_extension}

The baseline model and its uncorrelated-multiple-predictor extension (Section \ref{sec:lasso}) both assume that every candidate predictor is statistically unrelated to the dividend innovation $\eta_t$, so that any equilibrium belief and any selection episode is purely a finite-sample artifact: the CMCE is always $\beta^*=0$, and LASSO selects a predictor only when sampling noise pushes $\hat\beta^{\text{OLS}}$ outside the band $(-\lambda,\lambda)$. We now ask what changes when one predictor is genuinely informative about fundamentals. We show that such a predictor (i) acquires a nonzero equilibrium belief, so that selection reflects real (mis-specified but population-level) predictability rather than only noise, and (ii) is selected by LASSO with a probability that is increasing in the strength of its correlation with $\eta_t$, at least when the LASSO penalty is not set so high as to dominate the maximal attainable population signal.

For tractability we work with exactly two predictors, one ``clean'' (uncorrelated with fundamentals, as in the baseline) and one correlated with fundamentals.

\subsection{Setup}

\begin{assumption}[Bivariate Predictors with One Fundamentals-Correlated Regressor]\label{ass:corr_predictor}
Agents observe two predictors $(x_{1,t},x_{2,t})$, i.i.d.\ across $t$ and mutually uncorrelated within each $t$,
\begin{equation}\label{eq:bivariate_predictors}
\begin{pmatrix}x_{1,t}\\x_{2,t}\end{pmatrix} \sim \mathcal N\left(\begin{pmatrix}0\\0\end{pmatrix},\begin{pmatrix}\sigma_{x_1}^2 & 0\\ 0 & \sigma_{x_2}^2\end{pmatrix}\right),
\end{equation}
independent across $t$. The dividend innovation is generated by
\begin{equation}\label{eq:eta_corr}
\eta_t = \rho_{x\eta}\,\frac{\sigma_d}{\sigma_{x_2}}\,x_{2,t-1} + \sigma_d\sqrt{1-\rho_{x\eta}^2}\,\tilde\eta_t, \qquad \tilde\eta_t\stackrel{iid}{\sim}\mathcal N(0,1),
\end{equation}
where $\rho_{x\eta}\in(-1,1)$ and $\tilde\eta_t$ is independent of $\{x_{1,s},x_{2,s}\}_{s=-\infty}^{\infty}$ at every lead and lag. The predictor $x_{1,t}$ is unaffected and remains independent of $\{\eta_s\}$ at every lead and lag, exactly as in the baseline model.
\end{assumption}

We use $\rho_{x\eta}$ rather than $\rho$ for this new correlation parameter because $\rho$ already denotes the correlation between the dividend and consumption shocks $(\eta^d_t,\eta^c_t)$ in Section \ref{sec:model}; reusing it here would create a notational collision between two economically distinct correlations.

The construction in \eqref{eq:eta_corr} deliberately correlates $\eta_t$ with $x_{2,t-1}$ rather than with the contemporaneous $x_{2,t}$. This lag matches the timing convention used throughout the paper: the RLS/LASSO learning recursions pair the return innovation $\eta_{t-1}$ entering $r_{t-1}$ with the regressor $x_{t-2}$ that is one period older (see the ALM \eqref{eq:alm_RLS} and Definition \ref{def:cmce}), and the finite-sample benchmark of Appendix \ref{app:proof_LASSO} pairs $x_i$ with $r_i\approx\eta_i$ using an index $i$ that already encodes this one-period offset. If $\eta_t$ were instead correlated with the \emph{contemporaneous} $x_{2,t}$, the correlation would never appear in any of the population moments that govern learning or selection, because every regression and every CMCE condition in this paper relates $r_{t-1}$ (and hence $\eta_{t-1}$) to $x_{t-2}$, never to $x_{t-1}$. The lag-1 structure in \eqref{eq:eta_corr} is therefore the only specification under which $\rho_{x\eta}$ has any effect on the learning dynamics; we adopt it for this reason rather than as an arbitrary modeling choice.

Three properties of \eqref{eq:eta_corr} are immediate and used repeatedly below:
\begin{enumerate}
\item $\operatorname{Var}(\eta_t)=\rho_{x\eta}^2\sigma_d^2+\sigma_d^2(1-\rho_{x\eta}^2)=\sigma_d^2$ for every $\rho_{x\eta}$, so the unconditional dividend-shock variance is unaffected by the correlation and all FIRE/no-arbitrage steps of Section \ref{sec:model} go through unchanged.
\item $\operatorname{Corr}(x_{2,t-1},\eta_t)=\rho_{x\eta}$ exactly (not just to leading order), by direct computation from \eqref{eq:eta_corr} and \eqref{eq:bivariate_predictors}.
\item $x_{1,t}\perp \eta_s$ for all $t,s$, and $\rho_{x\eta}=0$ collapses \eqref{eq:eta_corr} to the baseline i.i.d.\ specification of Section \ref{sec:model}, so the baseline model is nested as a special case.
\end{enumerate}

\subsection{Equilibrium Belief and Local Stability}

With two predictors the PLM is $r_{t+1}=x_{1,t}\beta_1+x_{2,t}\beta_2+u_{t+1}$, and the ALM generalizes \eqref{eq:alm_RLS} to
\begin{equation}\label{eq:alm_bivariate}
r_{t-1}=\eta_{t-1}+g_\beta(x_{1,t-2},x_{2,t-2})-g_\beta(x_{1,t-1},x_{2,t-1}), \qquad g_\beta(x_1,x_2):=\phi(x_1\beta_1+x_2\beta_2),
\end{equation}
where $\phi(z):=\log(1-\kappa e^{w\tanh(z/w)})$ is exactly the univariate function defining $g_\beta(x)=\phi(x\beta)$ in \eqref{eq:g_beta_def}, so that $\phi'(z)=g_1'(z)$ coincides with the derivative computed in the proof of Proposition \ref{prop:RLS}, and $T(\beta)=\beta\,\mathbb E[\phi'(x\beta)]$ for $x\sim\mathcal N(0,\sigma_x^2)$ is literally Proposition \ref{prop:RLS}'s map. The bivariate CMCE condition is the natural generalization of Definition \ref{def:cmce}: $\beta_j=T_j(\beta_1,\beta_2)$ for $j=1,2$, where $T_j(\beta):=\operatorname{Cov}(r_{t-1},x_{j,t-2})/\sigma_{x_j}^2$.

\begin{proposition}[Equilibrium Shift and Local Stability under a Fundamentals-Correlated Predictor]\label{prop:corr_equilibrium}
Under Assumption \ref{ass:corr_predictor}, the bivariate CMCE conditions reduce to $\beta_1^*=0$ and
\begin{equation}\label{eq:beta2_star_implicit}
\beta_2^* - T(\beta_2^*) = \rho_{x\eta}\,\frac{\sigma_d}{\sigma_{x_2}},
\end{equation}
where $T(\cdot)$ is the univariate map of Proposition \ref{prop:RLS}. There exists $\bar\rho\in(0,1]$ and a unique, continuously differentiable function $\beta_2^*:(-\bar\rho,\bar\rho)\to\mathbb R$ solving \eqref{eq:beta2_star_implicit} with $\beta_2^*(0)=0$, satisfying
\begin{equation}\label{eq:beta2_star_deriv}
\left.\frac{d\beta_2^*}{d\rho_{x\eta}}\right|_{\rho_{x\eta}=0} = (1-\kappa)\,\frac{\sigma_d}{\sigma_{x_2}} > 0,
\end{equation}
so that $\beta_2^*(\rho_{x\eta})$ has the same sign as $\rho_{x\eta}$ and $\beta_2^*(\rho_{x\eta})\approx(1-\kappa)(\sigma_d/\sigma_{x_2})\,\rho_{x\eta}$ for $\rho_{x\eta}$ near zero. For $\rho_{x\eta}\neq 0$ in this neighborhood, $\beta_2^*\neq0$: the predictor correlated with fundamentals carries a genuinely nonzero equilibrium belief, in contrast to $\beta_1^*=0$ for the uncorrelated predictor. The equilibrium $(\beta_1^*,\beta_2^*(\rho_{x\eta}))$ is locally asymptotically stable under RLS learning for every $\kappa\in(0,1)$ and every $\rho_{x\eta}$ in this neighborhood of $0$.
\end{proposition}
\textit{Proof.} In Appendix \ref{app:proof_RLS_corr}.

\begin{remark}[Scope of Proposition \ref{prop:corr_equilibrium}]\label{rem:scope_corr}
The proposition is stated as a local result near $\rho_{x\eta}=0$, obtained via the implicit function theorem, because the global injectivity of $\beta\mapsto\beta-T(\beta)$ used to deliver a unique $\beta_2^*$ for \emph{every} $\rho_{x\eta}\in(-1,1)$ is not established by the existing machinery in Appendix \ref{app:proof_RLS}: that appendix proves $\operatorname{sign}[T(\beta)]=-\operatorname{sign}(\beta)$ for all $\beta\neq0$, which guarantees $\beta-T(\beta)$ is never zero away from the origin, but does not by itself guarantee monotonicity (and hence injectivity) of $\beta-T(\beta)$ away from a neighborhood of $0$. In numerical evaluation of $T(\beta)$ across a wide range of $\kappa\in(0,1)$ and $\beta$, we do not find any violation of global monotonicity, so a stronger, global version of Proposition \ref{prop:corr_equilibrium} may well be true; we leave a rigorous global proof (e.g.\ via an explicit bound $\sup_\beta T'(\beta)<1$) for future work and report only the scope that the implicit function theorem rigorously delivers.
\end{remark}

\subsection{Finite-Sample Selection Probability}

We now turn to the finite-sample rolling-window benchmark of Appendix \ref{app:proof_LASSO} (the paragraph following the proof of Proposition \ref{prop:LASSO_estimate}), rather than the infinite-history CMCE, since this is the object that connects most directly to the LASSO selection mechanism. As in that benchmark, we work under the null that $r_i\approx\eta_i$ within the rolling window of length $L$, and we relabel time so that the index $i$ used for $x_{2,i}$ already represents the same one-period-earlier pairing used in Assumption \ref{ass:corr_predictor} and in the dynamic model's pairing of $x_{t-2}$ with $r_{t-1}$; with this relabeling, $\operatorname{Corr}(x_{2,i},\eta_i)=\rho_{x\eta}$ within the benchmark's own indexing, consistent with \eqref{eq:eta_corr}.

\begin{proposition}[Finite-Sample Selection Probability under a Fundamentals-Correlated Predictor]\label{prop:corr_selection}
Under Assumption \ref{ass:corr_predictor} and the finite-sample rolling-window benchmark of Appendix \ref{app:proof_LASSO} with window length $L\geq 1$, the OLS estimator of $\beta_2$ satisfies
\begin{equation}\label{eq:beta2hat_L}
\hat\beta_{2,L} = \rho_{x\eta}\,\frac{\sigma_d}{\sigma_{x_2}} + \sqrt{1-\rho_{x\eta}^2}\,\frac{\sigma_d}{\sigma_{x_2}\sqrt L}\,T_L,
\end{equation}
where $T_L$ is exactly Student-$t$ distributed with $L$ degrees of freedom, the same object appearing in Appendix \ref{app:proof_LASSO}'s baseline benchmark. Let $c:=\lambda\sigma_{x_2}/\sigma_d$ denote the LASSO threshold expressed relative to the maximal attainable population signal $\sigma_d/\sigma_{x_2}$ (the slope obtained as $\rho_{x\eta}\to1$), and let
\begin{equation}\label{eq:Prho_gaussian}
P(\rho_{x\eta}) := \mathbb P\bigl(|\hat\beta_{2,L}|>\lambda\bigr) \;\approx\; \left[1-\Phi\!\left(\sqrt L\,\frac{c-\rho_{x\eta}}{\sqrt{1-\rho_{x\eta}^2}}\right)\right] + \Phi\!\left(-\sqrt L\,\frac{c+\rho_{x\eta}}{\sqrt{1-\rho_{x\eta}^2}}\right)
\end{equation}
denote the selection probability under the large-$L$ Gaussian approximation to $T_L$ (the exact statement uses the noncentral generalization of Appendix \ref{app:proof_LASSO}'s Student-$t$ formula). If $c\leq 1$ (equivalently $\lambda\leq\sigma_d/\sigma_{x_2}$), then $P(\rho_{x\eta})$ is strictly increasing in $\rho_{x\eta}$ for all $\rho_{x\eta}\in(0,1)$.
\end{proposition}
\textit{Proof.} In Appendix \ref{app:proof_LASSO_corr}.

\begin{remark}[The Case $c>1$]\label{rem:c_greater_1}
If $\lambda>\sigma_d/\sigma_{x_2}$, even a noiseless, perfectly correlated predictor ($\rho_{x\eta}\to1$) would have population slope $\sigma_d/\sigma_{x_2}<\lambda$ and would never be selected on population grounds alone; the proof in Appendix \ref{app:proof_LASSO_corr} shows that in this case the exact $P(\rho_{x\eta})$ in \eqref{eq:Prho_gaussian} is not monotonic on $(0,1)$: it rises from $0$, peaks near $\rho_{x\eta}=1/c$, and falls back to $0$ as $\rho_{x\eta}\to1$, because finite-sample noise — the only channel through which a predictor with sub-threshold population slope can ever be selected — is mechanically eliminated as $\rho_{x\eta}\to1$. We do not view this as a defect of the model; rather, it shows that ``increasing correlation strength increases selection'' is the empirically relevant regime precisely when the threshold is not set above the maximal attainable population signal, which Proposition \ref{prop:corr_selection} makes precise via the explicit condition $c\leq1$.
\end{remark}
